\clearpage
\chapter*{2025 Febrero}
\addcontentsline{toc}{chapter}{2025 Febrero}

\begin{enunciado}
    \hypertarget{2025 B 1}{}
    \textbf{1. (2,5 puntos)} Demostrar que en un espacio topológico \((X,T)\) de Hausdorff, si \(M_1,M_2\) son dos subconjuntos compactos disjuntos, entonces existen abiertos \(U_1\) y \(U_2\) que los separan; es decir \(U_1 \supset M_1, U_2 \supset M_2\) y \(U_1 \cap U_2 = \emptyset\).
\end{enunciado}

\begin{enunciado}
    \hypertarget{2025 B 2}{}
    \textbf{2. (3,5 puntos)} En el conjunto \(\mathbb{R}\) de los números reales, definimos la topología
    \[
        T = \{\mathbb{R}\} \cup P(\mathbb{R} - M), \text{ partes de } \mathbb{R} - M, \text{ siendo } M = (1,3)
    \]
    \begin{enumerate}[label=\alph*)]
        \item Estudiar si \((\mathbb{R},T)\) verifica el I axioma de numerabilidad
        \item Estudiar si \((\mathbb{R},T)\) es compacto
    \end{enumerate}
\end{enunciado}

\begin{enunciado}
    \hypertarget{2025 B 3}{}
    \textbf{3. (4 puntos)}
    \begin{enumerate}[label=\alph*)]
        \item Probar que un espacio topológico \((X,T)\) es discreto si y solamente si toda aplicación \(f : (X,T) \to (\mathbb{R},T_u)\) es continua.
        \item Sea el espacio \((X,T)\), los conjuntos \(Y,Y^*\) y las aplicaciones \(f : (X,T) \to Y\), \(g : Y \to Y^*\). Demostrar que si \(S\) es la topología final para \(f\), entonces \(S^*\) es la topología final para \(g : (Y,S) \to Y^*\) si y sólo si, \(S^*\) es la topología final para \(g \circ f : (X,T) \to Y^*\).
    \end{enumerate}
\end{enunciado}

